\section{Setting, objectives and metrics}
\label{sec:setting}

\paragraph{Data and units.}
Item $x_i$ has human distribution $p_i$ over ratings $r\in\{0,\dots,K-1\}$ and mean $t_i=\sum_r r\,p_i(r)$.
A judge forecast $q_i$ has mean $m_i=\sum_r r\,q_i(r)$.
A training budget $N$ is the total number of individual human ratings spent on fitting.
It can sit on $N$ items with one rating each ($d=1$ rating per item) or on $N/d$ items with $d$ ratings each.
Test scores average over held-out items, and each held-out item is scored against all of its collected ratings.

\paragraph{Objectives.}
Pearson fitting maximizes the correlation between forecast means and the mean training rating of each item.
Cross-entropy (CE) fitting minimizes $-\log q_i(r_{ij})$ averaged over individual training ratings.
RPS fitting and mean squared error (MSE) fitting minimize the squared cumulative-probability gap and the squared error $(m_i-r_{ij})^2$ on the same ratings (\appref{app:protocols}).
For a frozen judge, a fit chooses among 27 calibration settings of the judge's own probabilities $p$,
\begin{equation}
q\;\propto\;p^{\,t}\exp\!\bigl(b\,u_k+c\,u_k^{2}\bigr),
\label{eq:head}
\end{equation}
with $t\in\{\tfrac12,1,2\}$, $b,c\in\{-1,0,1\}$ and $u_k$ the centered, scaled category index.
Search is exhaustive; ties go to the setting nearest the identity; every objective sees the same settings and the same training ratings.
For a trained judge, LoRA weights~\citep{hu2022lora} and a probability read-out are updated by each objective under one optimizer, one schedule and one validation rule.

\paragraph{Metrics.}
The primary metric is the normalized individual-rating ranked probability score,
\begin{equation}
\RPS(q,r)=\frac{1}{K-1}\sum_{k=0}^{K-2}\bigl(F_q(k)-\Ind\{r\le k\}\bigr)^2 ,
\label{eq:rps}
\end{equation}
where $F_q$ is the cumulative distribution of $q$.
It lies in $[0,1]$, lower is better, and it is strictly proper~\citep{epstein1969scoring,gneiting2007strictly}.
We pair it with CE in nats (lower is better) and with total variation distance $\TVD=\tfrac12\sum_r|h(r)-q(r)|$ to each item's empirical human distribution $h$, averaged equally over items.
TVD is a diagnostic of the whole distribution, not a proper score for individual ratings.
Every table lists reference forecasts: a uniform forecast, the judge's raw output, and where relevant the training histogram.
Pearson, Spearman and Kendall correlations are secondary diagnostics of ranking.

\paragraph{Differences, uncertainty and break-even.}
A difference always subtracts the same metric on the same data for two named procedures.
We write it in words, for example ``test RPS of the Pearson-fitted judge minus test RPS of the CE-fitted judge'', and positive values favor CE.
Intervals are bootstraps over held-out items (over prompts for HelpSteer2, over folds with an overlap-adjusted $t$ interval for the homework studies), with simultaneous max-$t$ bands across budgets where stated; draws and levels are listed per experiment in \appref{app:protocols}.
The break-even budget $N^\ast$ is the smallest budget from which the lower bound for ``Pearson-fit minus CE-fit'' and for ``uniform minus CE-fit'' both stay above zero at every larger budget tested.
Selection of a fitted judge is by training data (and validation data where stated), never by the test ratings.
